\paragraph{证书化意义：有限核、谱裕度与离线验证器}
有限模态 Hardy 投影的核表示把连续对象的评估约化为有限维数据与显式有理核的积分（\S\ref{subsec:cdim-hardy-finite-mode-projection}）。
因此，在“有限数据证书 + 确定性离线验证器”的制度中，关键在于为数值实现提供可审计误差界与谱裕度封口；该语法在 Toeplitz--PSD 证书侧同样出现。
\par
具体地，设 \(\mathbf{T}_N=(c_{j-k})_{0\le j,k\le N}\) 为 Hermitian Toeplitz 主子矩阵，\(\widehat{\mathbf{T}}_N=(\widehat c_{j-k})\) 为其可审计近似。
若存在谱裕度 \(\lambda_{\min}(\widehat{\mathbf{T}}_N)\ge \eta>0\) 且算子误差满足
\[
\bigl\|\widehat{\mathbf{T}}_N-\mathbf{T}_N\bigr\|_{\mathrm{op}}\le \eta/2,
\]
则可推出 \(\mathbf{T}_N\succeq 0\)（定理 \ref{thm:xi-spectral-margin-finite-psd-certificate}）。
此外，该定理给出逐项系数误差到算子范数误差的显式转写：若 \(\max_{0\le n\le N}|\widehat c_n-c_n|\le \varepsilon\)，则
\[
\bigl\|\widehat{\mathbf{T}}_N-\mathbf{T}_N\bigr\|_{\mathrm{op}}
\le \sum_{k=-N}^{N}|\widehat c_k-c_k|
\le (2N+1)\varepsilon.
\]
因此，“有限核—误差界—谱裕度”可被视为跨接口的统一证书语法：输入为有限维数据与可核验界，输出为严格半正定性或禁入结论，并可被完全离线复核。

\endinput

